\subsection{Sprint 1}


Para a Sprint 1, nosso cronograma previa o amadurecimento das definições técnicas e o início da estruturação do projeto através do refinamento do protótipo e criação dos repositórios. O período foi marcado pela participação ativa da equipe em encontros com os \textit{stakeholders} e retrospectivas internas, visando garantir o entendimento das verdadeiras necessidades de negócio do SalesPilot para pavimentar o início das atividades de desenvolvimento e sua priorização de modo a gerar o maior valor para o cliente final.

Desse modo, ao longo dessa etapa, como \ac{ages} III, assumi a responsabilidade integral pela arquitetura do frontend, realizando a concepção completa do projeto de forma individual. Como outros colegas estavam designados a tarefas diferentes (como a arquitetura do backend e concepção da Wiki do projeto), em vez de um debate coletivo, realizei uma apresentação detalhada desta proposta de arquitetura e da \textit{technology stack} escolhida em um momento síncrono dedicado ao longo da semana. A estrutura apresentada foi fundamentada no uso de React 19 \cite{React}, TypeScript \cite{Typescript} e Vite \cite{Vite}, integrando o TanStack Router \cite{TanStackRouter} para um roteamento tipado e seguro, e o Zustand \cite{Zustand} para um gerenciamento de estado leve, eficiente e \textit{beginner-friendly}. Para o consumo da API e gerenciamento de cache, selecionei o React Query \cite{TanStackQuery} e optei por utilizar Zod \cite{Zod} para garantir uma validação rigorosa de schemas e dados em toda a aplicação. A interface segue os padrões da biblioteca de componentes Material UI (MUI) \cite{MUI}, como já havia sido definido na sprint anterior, e a confiabilidade do código é assegurada por testes unitários e de componentes utilizando Vitest \cite{Vitest} e React Testing Library \cite{ReactTestingLibrary}.

Dediquei também um tempo considerável à estrutura de pastas do projeto visando estabelecer uma organização lógica e intuitiva que facilitasse a navegação e a rápida familiarização de todos os integrantes com a base de código, principalmente dos colegas AGES I e II. Para isso me baseei em dois principais padrões de arquitetura: \textit{feature/domain-based architecture} \cite{FeatureBasedArchitecture} (agrupamento baseado em domínios) e \textit{smart-dumb} ou \textit{presentational-container components} \cite{ContainerPresentationalPattern} (separação de responsabilidades entre lógica e renderização). A escolha desses padrões visa promovar uma visibilidade maior dos testes além do repositório como um todo e simplificar significativamente o processo de revisão de código, garantindo diversos benefícios como a alta escalabilidade do sistema, a facilitação da reutilização de componentes em diferentes partes da aplicação, a redução da carga cognitiva para o time e uma manutenção muito mais ágil, assegurando que o projeto permaneça modular e escalável independentemente do aumento de complexibilidade do projeto a longo prazo.

Apesar do sucesso na apresentação do protótipo aos stakeholders e do feedback positivo recebido, a retrospectiva ao final da sprint demonstrou alguns problemas organizacionais e de responsabilidade de alguns colegas para com o time. Enfrentamos atrasos recorrentes e ausências constantes de alguns integrantes nos encontros presenciais e durante o desenvolvimento, o que prejudicou o ritmo de trabalho. É bastante desmotivador perceber como o descaso de certos membros com o projeto gera uma sobrecarga injusta sobre aqueles que estão de fato comprometidos. Além disso, esse desequilíbrio também acarretou problemas de gerenciamento de tempo, fazendo com que a sprint terminasse de forma apressada, colocando em risco a qualidade final das entregas. 

Essas adversidades reforçaram a importância de uma comunicação clara e da necessidade de uma postura mais assertiva na gestão de pessoas. Ficou evidente que o sucesso do SalesPilot depende de um esforço coletivo de todos como uma só equipe. Como próximos passos será necessário reforçar a política de presença, rever a alocação dos integrantes e estabelecer prazos para os objetivos, visando garantir que as tarefas sejam bem distribuídas, com maior visibilidade e que o comprometimento e a clareza para com o projeto seja uma regra reforçada e cultivada entre todos os membros da equipe.